\sectionkicker{Source-backed technical notes}
\subsection*{Open weightsから「再現できるモデルスタック」へ: Technical Notes}
本欄は記事本文で圧縮した一次資料上の情報を、比較・再検証しやすい形で整理したものである。時系列、確認済みの主張、留意点、一次資料URLを掲載する。
\smallskip
\noindent{\small\textit{Technical Notesでは、一次資料から確認した公開・提供・研究上の事実を記載する。数値・能力評価や提供元・著者による比較表現は、独立再現済みの結果ではなく帰属claimとして扱う。}}\par
\medskip
\subsection*{Theme at a glance}
\begin{center}
\footnotesize
\begin{tabularx}{\linewidth}{@{}p{0.20\linewidth}p{0.10\linewidth}p{0.14\linewidth}X@{}}
\toprule
資料 & 位置づけ & 種別 & 時系列 \\
\midrule
Meta Llama 3 & 主要資料 & オープンウェイト & 2024-04-18             \\
OLMo / OLMo 1.7 & 主要資料 & オープンウェイト & 2024-02-01, 2024-04-17             \\
DeepSeek-V2 / DeepSeek-Coder-V2 & 主要資料 & オープンウェイト & 2024-05-06, 2024-06-17             \\
DBRX & 補足資料 & オープンウェイト & 2024-03-27             \\
Mixtral 8x22B & 補足資料 & オープンウェイト & 2024-04-17             \\
Qwen2 & 補足資料 & オープンウェイト & 2024-06-07             \\
Gemma family (Gemma / PaliGemma / Gemma 2) & 補足資料 & オープンウェイト & 2024-02-21, 2024-05-14, 2024-06-27             \\
OpenELM & 補足資料 & オープンウェイト & 2024-04-22             \\
\bottomrule
\end{tabularx}
\end{center}
\normalsize

\Needspace{0.38\textheight}
\begin{technicalnote}{Meta Llama 3}{主要資料}
\begingroup
% reader-facing Technical Notes break policy
\widowpenalty=10000
\clubpenalty=10000
\displaywidowpenalty=10000
\begin{tabularx}{\linewidth}{@{}>{\bfseries}p{0.22\linewidth}X@{}}
組織 & Meta \\
種別 & オープンウェイト \\
時系列 & 2024-04-18 \\
\end{tabularx}
\smallskip
\begin{minipage}{\linewidth}
% reader-facing Technical Notes event-only fact/source tail group
{\bfseries 一次資料から整理したtechnical points}
\begin{itemize}[leftmargin=1.5em,itemsep=0.35em]
\item \textbf{一次情報で確認できる事実}: Metaの一次資料により、Meta Llama 3について2024-04-18にモデル公開を確認した。 Meta Llama 3の2024年4月公開分は8B/70Bのpretrained版とinstruction-fine-tuned版で、decoder-only Transformerを基盤に128K-token vocabularyのtokenizer、両サイズでgrouped query attention (GQA)、8,192-token sequence trainingを採用した。Metaはpretraining corpusを公開情報由来15T token超と説明しており、benchmark性能はMetaによる評価として扱う。
\end{itemize}
\begin{minipage}{\linewidth}
% reader-facing Technical Notes generic boundary/source tail group
\begin{samepage}
{\bfseries 一次資料}
\begingroup\sloppy
\Urlmuskip=0mu plus 2mu\relax
\begin{itemize}[leftmargin=1.5em,itemsep=0.25em]
\item {\scriptsize\url{https://ai.meta.com/blog/meta-llama-3/}}
\end{itemize}
\endgroup
\end{samepage}
\end{minipage}
\end{minipage}
\endgroup
\end{technicalnote}
\medskip

\Needspace{0.38\textheight}
\begin{technicalnote}{OLMo / OLMo 1.7}{主要資料}
\begingroup
% reader-facing Technical Notes break policy
\widowpenalty=10000
\clubpenalty=10000
\displaywidowpenalty=10000
\begin{tabularx}{\linewidth}{@{}>{\bfseries}p{0.22\linewidth}X@{}}
組織 & Ai2 \\
種別 & オープンウェイト \\
時系列 & 2024-02-01, 2024-04-17 \\
\end{tabularx}
\smallskip
\begin{minipage}{\linewidth}
% reader-facing Technical Notes event-only fact/source tail group
{\bfseries 一次資料から整理したtechnical points}
\begin{itemize}[leftmargin=1.5em,itemsep=0.35em]
\item \textbf{一次情報で確認できる事実}: Ai2の一次資料により、OLMo / OLMo 1.7について2024-02-01にモデル公開、2024-04-17にモデル更新を確認した。 対象event近傍の一次資料から 7B parameter scale を確認できる。
\end{itemize}
\begin{minipage}{\linewidth}
% reader-facing Technical Notes generic boundary/source tail group
\begin{samepage}
{\bfseries 一次資料}
\begingroup\sloppy
\Urlmuskip=0mu plus 2mu\relax
\begin{itemize}[leftmargin=1.5em,itemsep=0.25em]
\item {\scriptsize\url{https://allenai.org/blog/olmo-1-7-7b-a-24-point-improvement-on-mmlu-92b43f7d269d}}
\item {\scriptsize\url{https://allenai.org/blog/olmo-open-language-model-87ccfc95f580}}
\end{itemize}
\endgroup
\end{samepage}
\end{minipage}
\end{minipage}
\endgroup
\end{technicalnote}
\medskip

\Needspace{0.38\textheight}
\begin{technicalnote}{DeepSeek-V2 / DeepSeek-Coder-V2}{主要資料}
\begingroup
% reader-facing Technical Notes break policy
\widowpenalty=10000
\clubpenalty=10000
\displaywidowpenalty=10000
\begin{tabularx}{\linewidth}{@{}>{\bfseries}p{0.22\linewidth}X@{}}
組織 & DeepSeek \\
種別 & オープンウェイト \\
時系列 & 2024-05-06, 2024-06-17 \\
\end{tabularx}
\smallskip
\begin{minipage}{\linewidth}
% reader-facing Technical Notes event-only fact/source tail group
{\bfseries 一次資料から整理したtechnical points}
\begin{itemize}[leftmargin=1.5em,itemsep=0.35em]
\item \textbf{一次情報で確認できる事実}: DeepSeekの一次資料により、DeepSeek-V2 / DeepSeek-Coder-V2について2024-05-06にモデル公開、2024-06-17にモデル公開を確認した。 DeepSeek-V2自身はMixture-of-Experts構成とMulti-head Latent Attentionを用い、KV-cache量を抑える設計として説明される。DeepSeek-Coder-V2はDeepSeek-V2系のMoE code modelとして公開され、128K contextをサポートすると説明される。
\end{itemize}
\begin{minipage}{\linewidth}
% reader-facing Technical Notes generic boundary/source tail group
\begin{samepage}
{\bfseries 一次資料}
\begingroup\sloppy
\Urlmuskip=0mu plus 2mu\relax
\begin{itemize}[leftmargin=1.5em,itemsep=0.25em]
\item {\scriptsize\url{https://github.com/deepseek-ai/DeepSeek-Coder-V2}}
\item {\scriptsize\url{https://github.com/deepseek-ai/DeepSeek-V2}}
\end{itemize}
\endgroup
\end{samepage}
\end{minipage}
\end{minipage}
\endgroup
\end{technicalnote}
\medskip

\Needspace{0.38\textheight}
\begin{technicalnote}{DBRX}{補足資料}
\begingroup
% reader-facing Technical Notes break policy
\widowpenalty=10000
\clubpenalty=10000
\displaywidowpenalty=10000
\begin{tabularx}{\linewidth}{@{}>{\bfseries}p{0.22\linewidth}X@{}}
組織 & Databricks \\
種別 & オープンウェイト \\
時系列 & 2024-03-27 \\
\end{tabularx}
\smallskip
\begin{minipage}{\linewidth}
% reader-facing Technical Notes event-only fact/source tail group
{\bfseries 一次資料から整理したtechnical points}
\begin{itemize}[leftmargin=1.5em,itemsep=0.35em]
\item \textbf{一次情報で確認できる事実}: Databricksの一次資料により、DBRXについて2024-03-27にモデル公開を確認した。 対象event近傍の一次資料から 32K context / 12T training tokens / Mixture-of-Experts (MoE) を確認できる。
\end{itemize}
\begin{minipage}{\linewidth}
% reader-facing Technical Notes generic boundary/source tail group
\begin{samepage}
{\bfseries 一次資料}
\begingroup\sloppy
\Urlmuskip=0mu plus 2mu\relax
\begin{itemize}[leftmargin=1.5em,itemsep=0.25em]
\item {\scriptsize\url{https://www.databricks.com/blog/introducing-dbrx-new-state-art-open-llm}}
\end{itemize}
\endgroup
\end{samepage}
\end{minipage}
\end{minipage}
\endgroup
\end{technicalnote}
\medskip

\Needspace{0.38\textheight}
\begin{technicalnote}{Mixtral 8x22B}{補足資料}
\begingroup
% reader-facing Technical Notes break policy
\widowpenalty=10000
\clubpenalty=10000
\displaywidowpenalty=10000
\begin{tabularx}{\linewidth}{@{}>{\bfseries}p{0.22\linewidth}X@{}}
組織 & Mistral AI \\
種別 & オープンウェイト \\
時系列 & 2024-04-17 \\
\end{tabularx}
\smallskip
\begin{minipage}{\linewidth}
% reader-facing Technical Notes event-only fact/source tail group
{\bfseries 一次資料から整理したtechnical points}
\begin{itemize}[leftmargin=1.5em,itemsep=0.35em]
\item \textbf{一次情報で確認できる事実}: Mistral AIの一次資料により、Mixtral 8x22Bについて2024-04-17にモデル公開を確認した。 Mixtral 8x22Bはsparse Mixture-of-Expertsモデルで、総141B parametersのうち推論時のactive parametersは39BとMistral AIが説明している。64K-token context、native function calling、Apache 2.0でのweights公開が同じリリースの主要仕様であり、性能比較はMistral AIによる評価として扱う。
\end{itemize}
\begin{minipage}{\linewidth}
% reader-facing Technical Notes generic boundary/source tail group
\begin{samepage}
{\bfseries 一次資料}
\begingroup\sloppy
\Urlmuskip=0mu plus 2mu\relax
\begin{itemize}[leftmargin=1.5em,itemsep=0.25em]
\item {\scriptsize\url{https://mistral.ai/news/mixtral-8x22b}}
\end{itemize}
\endgroup
\end{samepage}
\end{minipage}
\end{minipage}
\endgroup
\end{technicalnote}
\medskip

\Needspace{0.38\textheight}
\begin{technicalnote}{Qwen2}{補足資料}
\begingroup
% reader-facing Technical Notes break policy
\widowpenalty=10000
\clubpenalty=10000
\displaywidowpenalty=10000
\begin{tabularx}{\linewidth}{@{}>{\bfseries}p{0.22\linewidth}X@{}}
組織 & Alibaba / Qwen \\
種別 & オープンウェイト \\
時系列 & 2024-06-07 \\
\end{tabularx}
\smallskip
\begin{minipage}{\linewidth}
% reader-facing Technical Notes event-only fact/source tail group
{\bfseries 一次資料から整理したtechnical points}
\begin{itemize}[leftmargin=1.5em,itemsep=0.35em]
\item \textbf{一次情報で確認できる事実}: Alibaba / Qwenの一次資料により、Qwen2について2024-06-07にモデル公開を確認した。 Qwen2は0.5B、1.5B、7B、57B-A14B、72Bのbase/instruction-tuned familyとして公開され、全サイズでGQAを採用した。Qwen Teamは英語・中国語に加えて27言語の学習データを含むこと、instruction-tuned 7B/72BでYARNを用いた最大128K context対応を説明しており、benchmark値は提供元評価として扱う。
\end{itemize}
\begin{minipage}{\linewidth}
% reader-facing Technical Notes generic boundary/source tail group
\begin{samepage}
{\bfseries 一次資料}
\begingroup\sloppy
\Urlmuskip=0mu plus 2mu\relax
\begin{itemize}[leftmargin=1.5em,itemsep=0.25em]
\item {\scriptsize\url{https://qwenlm.github.io/blog/qwen2/}}
\end{itemize}
\endgroup
\end{samepage}
\end{minipage}
\end{minipage}
\endgroup
\end{technicalnote}
\medskip

\Needspace{0.38\textheight}
\begin{technicalnote}{Gemma family (Gemma / PaliGemma / Gemma 2)}{補足資料}
\begingroup
% reader-facing Technical Notes break policy
\widowpenalty=10000
\clubpenalty=10000
\displaywidowpenalty=10000
\begin{tabularx}{\linewidth}{@{}>{\bfseries}p{0.22\linewidth}X@{}}
組織 & Google \\
種別 & オープンウェイト \\
時系列 & 2024-02-21, 2024-05-14, 2024-06-27 \\
\end{tabularx}
\smallskip
\begin{minipage}{\linewidth}
% reader-facing Technical Notes event-only fact/source tail group
{\bfseries 一次資料から整理したtechnical points}
\begin{itemize}[leftmargin=1.5em,itemsep=0.35em]
\item \textbf{一次情報で確認できる事実}: Googleの一次資料により、Gemma family (Gemma / PaliGemma / Gemma 2)について2024-02-21にモデル公開、2024-05-14にモデル公開、2024-06-27にモデル公開を確認した。 対象event近傍の一次資料から 2B parameter scale / 7B parameter scale / 9B parameter scale / 27B parameter scale を確認できる。
\end{itemize}
\begin{minipage}{\linewidth}
% reader-facing Technical Notes generic boundary/source tail group
\begin{samepage}
{\bfseries 一次資料}
\begingroup\sloppy
\Urlmuskip=0mu plus 2mu\relax
\begin{itemize}[leftmargin=1.5em,itemsep=0.25em]
\item {\scriptsize\url{https://blog.google/innovation-and-ai/technology/developers-tools/gemma-open-models/}}
\item {\scriptsize\url{https://blog.google/innovation-and-ai/technology/developers-tools/google-gemma-2/}}
\item {\scriptsize\url{https://developers.googleblog.com/en/gemma-family-and-toolkit-expansion-io-2024/}}
\end{itemize}
\endgroup
\end{samepage}
\end{minipage}
\end{minipage}
\endgroup
\end{technicalnote}
\medskip

\Needspace{0.38\textheight}
\begin{technicalnote}{OpenELM}{補足資料}
\begingroup
% reader-facing Technical Notes break policy
\widowpenalty=10000
\clubpenalty=10000
\displaywidowpenalty=10000
\begin{tabularx}{\linewidth}{@{}>{\bfseries}p{0.22\linewidth}X@{}}
組織 & Apple \\
種別 & オープンウェイト \\
時系列 & 2024-04-22 \\
\end{tabularx}
\smallskip
\begin{minipage}{\linewidth}
% reader-facing Technical Notes event-only fact/source tail group
{\bfseries 一次資料から整理したtechnical points}
\begin{itemize}[leftmargin=1.5em,itemsep=0.35em]
\item \textbf{一次情報で確認できる事実}: Appleの一次資料により、OpenELMについて2024-04-22にPaper First Submission、2024-04-22に研究公開を確認した。 OpenELMはopen language model familyで、Transformer各層へparameter budgetを一様配分せずlayer-wise scaling strategyで割り当てる設計を中心に提示する。Appleはtraining/evaluation frameworkに加え、public checkpoints・configurations・logsとApple device向けMLX conversion codeを公開している。
\end{itemize}
\begin{minipage}{\linewidth}
% reader-facing Technical Notes generic boundary/source tail group
\begin{samepage}
{\bfseries 一次資料}
\begingroup\sloppy
\Urlmuskip=0mu plus 2mu\relax
\begin{itemize}[leftmargin=1.5em,itemsep=0.25em]
\item {\scriptsize\url{http://arxiv.org/abs/2404.14619v2}}
\item {\scriptsize\url{https://machinelearning.apple.com/research/openelm}}
\end{itemize}
\endgroup
\end{samepage}
\end{minipage}
\end{minipage}
\endgroup
\end{technicalnote}
\medskip
